\section{Small two-power characters at coefficient seven on the same window}
\label{sec:coefficient-seven-actual-window}

\paragraph{Proved arithmetic input and exact scope.}
The new input is the original-coefficient noncritical strict theorem
AC-G197, applied precisely at (R7.3) and (R7.8), with its four
completed class-quotient certificates and full original-local proof.
The unused real-$168$ computation supplies no hypothesis.
This is a conditional-interface proof; the arithmetic and whole-proof
checks are recorded in the companion sources and \texttt{CHECK\_EVENTS.tsv}.

\paragraph{The exact original remainder.}
Keep every original datum, page, prescribed pair, completion, lattice
and local condition of AC-G196-window. For integers $h\geq2$,
$T\geq1$, let $Z_u,J_{u,0},J_{u,\rm mix},V_u,Y_u$ be exactly
the preceding sets in (CAW.1)--(CAW.4) and (RMW.2)--(RMW.3).
Thus $Z_u$ is the full old G190 omitted set; $V_u$ is the two
coefficient-three noncritical rows; $Y_u$ is the disjoint G195
split-prime, seven-prime and auxiliary residual-matching union.
At source size $u$ the unchanged actual host and weight range are
\[
 \begin{gathered}
 H_u=\mathbb Q(\zeta_{2^{h+2}},\zeta_7,\sqrt{-q},\sqrt{-d},
                          \omega,\sqrt[3]{269}),\quad K=\mathbb Q(\omega),\\
 u^{2+8\varepsilon}\leq2^h\leq2u^{2+8\varepsilon},\quad
 T=\lfloor u^{399}\rfloor,\quad0<\varepsilon<1/256.
 \end{gathered}\tag{S7W.1}
\]
All coefficient values and residue reductions use the prescribed
original $O_7$. Let $\tau_7$ denote the Teichm\"uller lift of
$\overline c_7$ as in (RQM.1), and put
\[
 \begin{split}
 C_7(u)={}&\{(7,(\delta\tau_7^j)|_{G_K},e):
       \delta\in\{1,\chi_{-4},\chi_{-8}\},\ 0\leq j<6,\ 1\leq e<T\}\\
 &\sqcup\{(7,(\chi_8\tau_7^j)|_{G_K},e):
                    0\leq j<6,\ 1\leq e<T,\ j-e\text{ odd}\},\\
 Y_7(u)={}&\{(7,\tau_7^e|_{G_K},e):1\leq e<T\},\qquad
 U_7(u)=C_7(u)\setminus Y_7(u).
 \end{split}\tag{S7W.2}
\]
Every undisplayed Kummer cubic and auxiliary factor is trivial.
If $p(\delta)\in\{0,1\}$ is defined by
$\delta(-1)=(-1)^{p(\delta)}$, set
\[
 b\in\{0,1\},\quad b\equiv p(\delta)+j-e\pmod2,\qquad
 \psi^{\rm nc}_{\delta,j,e}=\delta\tau_7^j\chi_{-3}^{\,b}.
 \tag{S7W.3}
\]
Its parity is $(-1)^e$, so it is the unique noncritical finite
rational extension used in the preceding remainder.

The exact G197 conclusion (R7.3) gives
\[
 \mathcal S^{\mathbb Q}_7\bigl((E_7/O_7)
                (\psi^{\rm nc}_{\delta,j,e}c_7^{-e})\bigr)=0
                  \qquad((7,\delta\tau_7^j,e)\in C_7(u)),
 \tag{S7W.4}
\]
on this original ring, with full localization at seven, the original
away-inertia kernels and the rational real-place convention.
No corresponding zero for the critical partner is asserted.

Define the explicit remaining sum
\[
 \begin{split}
 \mathcal R_{399}^{\rm match,7}(u)={}&
 \sum_{(\ell,\eta,e)\in J_{u,\rm mix}}
                       \frac{\log\#\mathcal S_{\ell,\eta}(e)}{b_\ell}\\
 &+\sum_{(\ell,\eta,e)\in J_{u,0}\setminus(V_u\cup Y_u\cup C_7(u))}
           \frac{\log\#\mathcal S^{\mathbb Q}_{\ell,
                             \psi_{\eta,e}^{\rm nc}}(e)}{b_\ell}.
 \end{split}\tag{S7W.5}
\]

\begin{theorem}[AC-G198-window: twenty further coefficient-seven terms per weight]
\label{AC-G198-window}
On every original datum and weight range above, all the displayed
sets are actual original-character sets and
\[
 \begin{gathered}
 C_7(u)\cap(Z_u\cup V_u\cup Y_u)=Y_7(u),\qquad
 U_7(u)\subset J_{u,0}\setminus(V_u\cup Y_u),\\
 |C_7(u)|=21(T-1),\quad |Y_7(u)|=T-1,\quad |U_7(u)|=20(T-1),\\
 \mathcal R_{399}^{\rm match,7}(u)=\mathcal R_{399}^{\rm match}(u).
 \end{gathered}\tag{S7W.6}
\]
The last equality deletes zero rational noncritical terms, not
original-$K$ strict groups or positive derived indices.
All the old omitted sets and their proven scopes are retained.
The newly absent indexing triples number $20(T-1)$, and the
union $Y_u\cup C_7(u)$ has cardinality
\[
              21(T-1)+J_+(T)+A_q(T)+A_d(T).
 \tag{S7W.7}
\]
Here $J_+,A_q,A_d$ retain their exact definitions from (RMW.5),
including exclusion of seven from $J_+$ and possibly zero auxiliary
counts. The set of original triples not represented in (S7W.5)
is $Z_u\sqcup V_u\sqcup Y_u\sqcup U_7(u)$.

Retain exactly $\mathcal E_u,\mathcal L_u,\mathcal A_u,\mathcal P_u$
and constants from (CAW.5) and (RMW.7). Then
\[
 \begin{aligned}
 |I_{399}-\mathcal R_{399}^{\rm match,7}|&\leq
          \mathcal E_u+\mathcal L_u+\mathcal A_u+C_2\mathcal P_u,\\
 |I_{399}^{\rm all7,2}-\mathcal R_{399}^{\rm match,7}|&\leq
          \mathcal E_u+\mathcal L_u+C_2\mathcal P_u.
 \end{aligned}\tag{S7W.8}
\]
The one full local majorant, one torsion sum and one critical
aggregate are unchanged. No new family multiplier or second
critical bound is inserted. The same uniform envelope remains
$O_\varepsilon(u^{806+16\varepsilon}(\log u)^2)$.
The old and new explicit numerical remainder premises are equal
on the identical window, threshold and slack. The full original
index premise follows with the existing threshold increase and
half-slack argument on that same window.
\end{theorem}

\begin{proof}
\proofstep{1. Literal original-coefficient input and rational parity.}
The original host contains the conductor-eight and conductor-seven
cyclotomic factors. The four characters
$1,\chi_{-4},\chi_8,\chi_{-8}$ are precisely the four rational
characters of order dividing two and conductor dividing eight. Their values
are $\pm1$ in the original $O_7$; those of $\tau_7$ are the
sixth roots in $\mathbb Z_7$, already in that same ring.
The full original character order $3\cdot2^h$ is divisible by six.
No coefficient extension or new residue embedding is needed.

Equation (S7W.3) gives
$p(\delta)+j+b\equiv e\pmod2$, proving the exact noncritical
sign. Its residual full coefficient is
$\delta\tau_7^{j-e}\chi_{-3}^{b}$, with all values interpreted
in the original residue field. It is exactly the original coefficient
of (R7.2)--(R7.3), with the unique noncritical parity lift.
Apply (R7.3) to obtain (S7W.4), using the same full seven-adic group,
away inertia restrictions and localization-set inflation as
G191 and G196. At coefficient seven the index-two lattice
comparison has determinant two, a unit in $O_7$.
The original $K$-strict group can retain a nonzero critical
component; only its rational noncritical term is removed here.

\proofstep{2. Exact old-set intersections.}
All triples of $C_7(u)$ have coefficient seven and trivial cubic
and auxiliary factors. Hence they miss every old coefficient-two,
coefficient-three and coefficient-$269$ family and the set $V_u$.
The only coefficient-seven rows of the old $Z_u$ are the high-two
families in (OPW.3). The inert-trivial family has no coefficient
seven because $7\equiv1\pmod3$. For G175's threshold,
\[
 v_2(7^2-1)=4,\quad21\cdot7\cdot2^8=37632,\quad c_7=16.
 \tag{S7W.9}
\]
Every such old coefficient-seven character has primitive two-part
of conductor at least $2^{16}$, whereas all two-parts in $C_7$
have conductor dividing eight. The old high-two family is thus
disjoint from $C_7$, whether or not it occurs at the current $h$.

The G196 set $Y_u$ has coefficient-seven part exactly $Y_7(u)$:
$Y_+$ excludes seven, and $q,d>7$. Its one character at weight
$e$ is $\tau_7^e$, the $\delta=1$, $j\equiv e\pmod6$ row of
$C_7(u)$. There are no other possible intersections.
Conductor independence makes the four two-parts and six seven-parts
distinct over $\mathbb Q$; their restrictions remain distinct over
$K$ since a character of conductor prime to three cannot equal
$\chi_{-3}$. This proves the first line of (S7W.6).

\proofstep{3. Counts and exact zero-term deletion.}
At every fixed weight, each of $\delta=1,\chi_{-4},\chi_{-8}$
allows all six exponents $j$. For $\delta=\chi_8$, exactly three
of $j=0,\ldots,5$ have parity opposite to $e$.
There are therefore $18+3=21$ distinct rows per weight.
Exactly one is the already deleted $Y_7$ row, leaving twenty.
Multiplying by $T-1$ proves all counts in (S7W.6)--(S7W.7),
including the empty case $T=1$.
For $T\geq2$, the row with $\delta=\chi_{-4},j=0,e=1$
exhibits proper reduction of the previously retained indexing set.

Every removed group is zero by the exact input (S7W.4).
All sums are finite by the retained G191/G192 finiteness proof.
Deleting these zero normalized terms proves the last line of
(S7W.6), on the original $b_7$ and all original local conditions.
No claim of strict numerical decrease follows from deleting
zero terms. The three $\chi_8$ rows with $j-e$ even at each weight
are not covered by this argument and remain in the remainder.

\proofstep{4. No new cost and the identical common window.}
Substitute the exact remainder equality into (RMW.7) to get
(S7W.8). The critical partners of every row of $C_7\setminus Y_7$
were already in the one G191 critical aggregate over $J_{u,0}$.
Neither the local error, the torsion estimate nor the critical
estimate is charged a second time. Source degree $D_u=144\cdot2^h$
and abelian degree $d_u=96\cdot2^h$ remain distinct.

Fix the original datum, page and $0<\varepsilon<1/256$.
The exact (TC.3) parameters $J,X_0,s$, prescribed pairs and
window $X^{1/J}\leq W$, $W^2\leq X$ remain those of G196.
Pointwise equality of the explicit remainders preserves their
numerical premises with precisely the same window and slack.
For comparison with the full index sum choose the same uniform
$U_0$ making the first error of (S7W.8) at most
$(s/2)\exp(u^\varepsilon/2)$, raise the threshold to
$\max(X_0,U_0^J)$, and retain that exact $W$.
The full premise holds with slack $s/2$ on $[W,W^2]$.
No separate existential windows are intersected. Every source
rank $R.t$, coefficient, page, pair and finite simultaneous
condition is retained. The finite-etale original-coefficient
and strict completion/inertia dictionary remains that of G196.
\end{proof}

\paragraph{Unchanged open scope.}
All mixed cubic terms, unsupported noncritical rational terms,
and in particular the three even-$j-e$ $\chi_8$ rows remain.
The uniform numerical bound for the remaining sum, the mixed
theta/Selmer realization and H04.4 remain open. The proof removes
exactly the twenty previously retained noncritical terms per weight,
with no new question row or automatic movement credit.
